\documentclass[12pt]{article}

\usepackage[utf8]{inputenc}
\usepackage{latexsym,amsfonts,amssymb,amsthm,amsmath}
\usepackage{tikz}
\usepackage{stanli}
\usetikzlibrary{calc,through}
\usepackage{multicol}
\usepackage{geometry} 
\usepackage{enumitem}

\definecolor{darkred}{rgb}{0.6,0,0}

\setlength{\parindent}{0in}
\setlength{\oddsidemargin}{-0.25in}
\setlength{\textwidth}{7in}
\setlength{\textheight}{9.3in}
\setlength{\topmargin}{-1in}
\setlength{\headheight}{18pt}

\newcommand\blfootnote[1]{%
  \begingroup
  \renewcommand\thefootnote{}\footnote{#1}%
  \addtocounter{footnote}{-1}%
  \endgroup
}

\title{ME473 Dynamic Finite Element Analysis Mini-Project}
\author{Stefano Burzio}
\date{}

\begin{document}
\noindent ME473 - Dynamic finite element analysis of structures\hfill EPFL\\
Stefano Burzio \hfill 2025\\
\vspace{0.3cm}

\hspace*{\fill}\textbf{\Large Mini-project 2}\hspace*{\fill}
\begin{center}
  \textbf{\Large Design of a 2d truss crane}
\end{center}

\hrulefill
\vspace{0.3cm}
\newline
{Project organization:}
\vspace{0.2cm}
\newline
\begin{minipage}[t]{0.38\textwidth}
  \begin{itemize}[itemsep=0.005cm]
    \item Groups: 3 to 4 students
    \item 20\% of final grade
    \item Pdf report: maximum 10 pages
  \end{itemize}
\end{minipage}
\hfill
\begin{minipage}[t]{0.58\textwidth}
  \begin{itemize}[itemsep=0.005cm]
    \item Programming language: \texttt{MATLAB}
    \item Submission: November 28, 2025
    \item Total workload: 40h (approx. 10h per student)
  \end{itemize}
\end{minipage}
\vspace{0.3cm}

\hrulefill
\vspace{0.3cm}

Your task is to design a crane truss structure for three different materials and make sure each individual structure you created fulfils the detailed static and dynamic requirements. The crane must lift a vertical point load of $25$ kN applied at the tip of the jib (node $3$), which extends vertically $h=20$ meters above the ground level and horizontally $\ell=14$ meters from the base left support (node $2$). The two supports at the base (nodes $1$ and $2$) are fixed to the ground and are separated by a distance of $d=2$ meters.

\subsubsection*{Anchor points and load configuration}
The crane is fixed at the following supports on a horizontal plane:
\begin{multicols}{2}
  \begin{itemize}
    \item Base left: (0, 0)
    \item Base right: (2, 0)
  \end{itemize}
\end{multicols}

\begin{center}
  \begin{tikzpicture}[scale=0.8]
    % Crane nodes
    \coordinate (BaseLeft) at (0,0);
    \coordinate (BaseRight) at (2,0);
    \coordinate (MastTop) at (0,8);
    \coordinate (JibTip) at (8,8);
    \coordinate (Diagonal) at (4,4);
    % Distances members
    \draw[latex-latex,dashed,thin] (BaseLeft) -- (MastTop) node[midway,left] {\small $h=20$ m};
    \draw[latex-latex,dashed,thin] (MastTop) -- (JibTip) node[midway,above] {\small $\ell=14$ m};
    \draw[latex-latex,dashed,thin] (BaseLeft) -- (BaseRight) node[midway,above] {\small $d=2$ m};
    % Supports
    \support{1}{BaseLeft};
    \support{1}{BaseRight};

    % Load at jib tip
    \draw[thick,latex-] (JibTip) -- ++(0,1.5) node[right] {25 kN};
    % Nodes
    \node[left=1pt,darkred] at (BaseLeft) {1};
    \node[right=1pt,darkred] at (BaseRight) {2};
    \node[right=1pt,darkred] at (JibTip) {3};
    % Draw points
    \foreach \p in {BaseLeft, BaseRight, JibTip} {
        \draw[thick,fill=white] (\p) circle (3pt);
      }
  \end{tikzpicture}
\end{center}

\subsubsection*{Crane design requirements}
The crane must be modeled as a planar truss. All joints are idealized as frictionless pinned connections. You are required to perform both a static and a dynamic analysis for each of the three materials listed below, treating each case as a distinct crane design. A distinct truss design must be developed for each of these three materials.
\begin{center}
  \begin{tabular}{l|c|c|c}
    {Material}       & {E (GPa)} & {Density (kg/m$^3$)} & {Cross-Sectional area} \\
    \hline
    Structural steel & 200       & 7850                 & 15–50 cm$^2$           \\
    Aluminum alloy   & 69        & 2700                 & 20–70 cm$^2$           \\
    Engineered wood  & 12        & 600                  & 30–100 cm$^2$          \\
  \end{tabular}
\end{center}

\subsubsection*{Static analysis}
Each crane design must satisfy the following:

\begin{enumerate}
  \item The structure must remain fully elastic with displacements below \textbf{60 mm}.
  \item The total material volume must not exceed:
        \begin{itemize}
          \item Structural steel: \textbf{1 m$^3$},
          \item Aluminum alloy: \textbf{1.7 m$^3$},
          \item Engineered wood: \textbf{2 m$^3$}.
        \end{itemize}
\end{enumerate}

\subsubsection*{Dynamic analysis}
A separate modal analysis must be performed for each of the three trusses designs corresponding to the different materials. The following requirement must be satisfied:
\begin{enumerate}
  \item[3.] Compute the first three natural frequencies and mode shapes. The first natural frequency must \emph{not} fall within \textbf{0 to 5 Hz}.
\end{enumerate}

\emph{Bonus}: To ensure dynamic stability and avoid resonance with common excitation sources, you should aim for a first natural frequency \textbf{above 20 Hz}, where feasible.

\subsubsection*{Deliverables}
\begin{enumerate}
  \item \texttt{MATLAB:}
        \begin{multicols}{2}
          \begin{itemize}
            \item Nodes and elements definitions
            \item Stiffness and mass matrices assembly
            \item Static and modal analysis
            \item Plots and visualizations
          \end{itemize}
        \end{multicols}

  \item \textbf{Report:}
        \begin{itemize}
          \item Crane layout with labeled pins and bars,
          \item Summary of \texttt{MATLAB} results (displacements, stresses, frequencies and mode shapes),
          \item Discussion of your findings.
        \end{itemize}
\end{enumerate}

\end{document}